% !TEX root = ../metatime_monograph.tex
\chapter{Neutron Electric Dipole Moment: Constraint Audit}
\label{ch:neutron_edm}

% CITATION_CONTEXT_V10_9_BEGIN
\paragraph{Established context.}
The link between a QCD CP phase and the neutron electric dipole moment follows
the chiral and effective-field-theory literature
\cite{crewther1979,pospelov2005}.  The experimental constraint is represented by
the modern neutron-EDM measurement \cite{Abel2020nEDM}.
% CITATION_CONTEXT_V10_9_END

\claimstatus{D -- falsification audit of the fixed strong-CP assignment.}

\section{Fixed mapping}
Using Chapter \ref{ch:strong_cp},
\[
  \theta_{\rm QCD}=2.2208116435\times10^{-10},
  \qquad
  C=2.4\times10^{-16}\ e\,\mathrm{cm},
\]
and therefore
\[
  \boxed{d_n^{\rm model}=5.3299\times10^{-26}\ e\,\mathrm{cm}}.
\]

\section{Experimental comparison}
The bound used in this release is
\[
  |d_n|<1.8\times10^{-26}\ e\,\mathrm{cm}\qquad(90\%\ \mathrm{CL}).
\]
The model value is $2.96$ times the bound.  The current fixed implementation
therefore fails this constraint.  This result is not described as ``at the
limit'' and is not deferred to a future experiment.

\section{Allowed interpretation}
The failure applies to the conjunction of the exponent $14$, the chosen
$\kappa$, and the fixed coefficient $C$.  It does not by itself prove that every
possible information-geometric strong-CP construction is excluded.  However,
changing the exponent to $15$, adding a suppression factor, or selecting a new
coefficient after seeing the bound would constitute a new model version and
cannot be counted as confirmation of the frozen formula.

\section{Relation to axion models}
The minimal Metatime implementation contains no Peccei--Quinn axion.  Discovery
of an axion would contradict the minimal claim that no additional strong-CP
degree of freedom is required; it would not automatically falsify every other
mathematical component of TIR.  Conversely, the absence of an axion is not
positive evidence for the present nEDM formula.
